% Section 2.2.3 of Prism main.tex, synchronized on 6 October 2026.
% Paths below refer to the existing Prism project.
\subsubsection{Validation}

The construction of the Spanish AI exposure measure relies on the assumption that occupational AI adoption patterns are sufficiently similar across countries for the Anthropic exposure measure to provide an informative ranking of AI exposure in Spain. The observed exposure measure uses global Claude usage classified under the U.S. O*NET taxonomy, rather than usage restricted to the U.S. labor market \citep{massenkoff_mccrory_2026}. We assess the validity of the resulting exposure measure through three complementary exercises: external validation, distributional validation, and face validity.

Our first validation exercise examines whether occupational AI usage patterns are similar in Spain and globally. Figure~\ref{fig:claude_usage_spain_global} compares the distribution of Claude usage across broad occupational groups using the latest Anthropic Economic Index. The two distributions are remarkably similar. Computer and mathematical occupations account for the largest share of Claude usage in both distributions, with education, arts and media, sales, business occupations, and office support also accounting for substantial shares. Differences between Spain and the global average are generally small, suggesting that occupational AI usage follows similar patterns in Spain and globally. This evidence supports the external validity of applying the Anthropic exposure ranking to the Spanish labor market. A similar occupation-level imputation strategy is also adopted by recent European evidence on generative AI and labor markets, including \citet{facius_iacono_2026}.

\begin{figure}[H]
\centering
\caption{Distribution of Claude usage by major occupational group: Spain and the global average}
\label{fig:claude_usage_spain_global}
\begin{threeparttable}
\includegraphics[width=0.9\linewidth]{uploads/prism-uploads/figure_anthropic_country_soc_major_group_spain_global_may2026.png}
\vspace{0.15cm}
\begin{tablenotes}[flushleft]
\footnotesize
\item Notes: The figure reports the percentage of Claude usage accounted for by each major occupational group according to the Anthropic Economic Index. Spain refers to Claude interactions originating in Spain; the global average refers to pooled worldwide Claude usage, rather than an unweighted average across countries.\newline
Source: Anthropic Economic Index (May 2026).
\end{tablenotes}
\end{threeparttable}
\end{figure}

Our second validation exercise evaluates whether the semantic mapping preserves the overall distribution of AI exposure when transferring the measure from the O*NET occupational classification to Spanish CNO4d occupations. Figure~\ref{fig:exposure_distribution_onet_spain} compares the distribution of exposure scores for the original Anthropic/O*NET occupations and the mapped Spanish occupations. The mapped distribution closely reproduces both the concentration of occupations with very low AI exposure and the upper tail of highly exposed occupations observed in the original O*NET data. This close correspondence indicates that the semantic mapping preserves the overall distribution of AI exposure despite the differences between the two occupational classifications.

\begin{figure}[H]
\centering
\caption{Distribution of AI exposure before and after mapping}
\label{fig:exposure_distribution_onet_spain}
\begin{threeparttable}
\includegraphics[width=0.7\linewidth]{uploads/prism-uploads/figure_ai_exposure_distribution_onet_vs_spain.png}
\vspace{0.15cm}
\begin{tablenotes}[flushleft]
\footnotesize
\item Notes: The figure compares the distribution of the original Anthropic AI exposure scores for O*NET occupations with the distribution obtained after mapping the exposure measure to Spanish CNO4d occupations using the nearest-neighbor semantic matching procedure.
\end{tablenotes}
\end{threeparttable}
\end{figure}

Finally, we assess the face validity of the mapped exposure measure by examining the occupations with the highest and lowest exposure levels. Online Appendix Table~\ref{tab:top_bottom_exposure_comparison} reports the ten occupations with the highest and lowest AI exposure after mapping the Anthropic measure to Spanish CNO4d occupations. The ranking is economically intuitive. Computer programmers, software specialists, financial analysts, and data-entry clerks appear among the most exposed occupations, whereas occupations involving predominantly manual or physical tasks are concentrated among the least exposed. The close correspondence between the ranking and prior expectations provides additional reassurance that the semantic mapping preserves economically meaningful relationships across occupational classifications.
